\section{The frozen comparison: results, augmentation and precision}
\label{sec:augprec}

\subsection{The frozen comparison}
\label{sec:frozen-comparison}

Table~\ref{tab:frozen} gives the frozen comparison that Section~\ref{sec:results-frozen}
summarises. We report no frozen-against-unfrozen contrast except one.
Eight configuration fields separate the frozen and unfrozen CORAL runs: head learning rate,
weight decay, cached features against images, augmentation, precision, the weight schedule, the
weight, and target-covariance averaging. A difference between them would mix the encoder
setting with the protocol. The exception is FreeMatch, whose frozen and unfrozen runs share
every training setting except the encoder and the further steps of training
(Section~\ref{sec:results-frozen}).

% GENERATED FILE -- do not edit by hand.
% Rebuild:  python make_all.py (generator: make_main_tables_v2.py)
% Table 1 (rewrite v1): the frozen comparison at two decimals. Sources as T3_round1.tex.
\begin{table}[tbp]
\centering
\caption{\textbf{The frozen comparison.} Means over folds 0--4, the development fold included (with across-fold sd), of TPR at FPR${}=10^{-3}$, encoder frozen, 8000 steps, last checkpoint. Differences against Baseline A are paired by fold over the same five folds, with 95\% intervals from the class-conditional paired bootstrap ($B=1000$). DANN ran on the development fold only, with a different adapter and budget, and on the original test pool; every other row uses the cleaned pool.}
\label{tab:frozen}
\footnotesize
\setlength{\tabcolsep}{4pt}
\begin{tabular}{l>{\TRAG}p{3.6cm}rr}
\toprule
Method & Family & TPR@$10^{-3}$ & vs.\ Baseline A, 95\% CI \\
\midrule
Baseline A & source only & $0.20 \pm 0.04$ & --- \\
CORAL & marginal alignment (moment matching) & $0.22 \pm 0.03$ & $+0.02$ $[+0.0001, +0.05]$ \\
AdaMatch & self-training & $0.43 \pm 0.05$ & $+0.23$ $[+0.17, +0.28]$ \\
FreeMatch & self-training & $0.43 \pm 0.03$ & $+0.23$ $[+0.19, +0.30]$ \\
\midrule
DANN (fold 0) & marginal alignment (adversarial) & $0.19$ & --- \\
\bottomrule
\end{tabular}
\end{table}


\subsection{Augmentation and precision}
\label{sec:augprec-detail}

The arms of the frozen comparison were not trained under identical conditions. The
feature-space arms ran on cached features, without augmentation, in full precision; the
image-space arms ran with weak augmentation (Section~\ref{sec:sim-pipeline}) in mixed precision, because weak and strong views
are the mechanism of self-training. We measured that pair of variables on the frozen
source-only baseline on the development fold, where only the head's \augprecTrainableNs{}
parameters train.

We ran three of the four cells of the $2 \times 2$ (Table~\ref{tab:augprec2}). Cell~A, no
augmentation at full precision, is the configuration of the feature-space arms; cell~C, weak
augmentation at half precision, is that of the image-space arms; cell~B, weak augmentation at
full precision, sits between them. We did not run cell~D, no augmentation at half precision.
Cell~A reads \augprecCellA{} on this pool, one positive below the Baseline~A value of
\augprecBaselineFoldZero{}; the two matched at $\baseFoldZeroPrereg{}$ on the original pool. We
ran all three cells locally so that they share one venue; the same cell~C configuration on
Kaggle reads \augprecCellCKaggle{}, \augprecVenueLenses{} positives higher
(\augprecVenueDiff{}, interval \augprecVenueCI).

The total difference $C - A$ is \augprecTotal{}, 95\% interval \augprecTotalCI. Its upper bound
sits one positive (\augprecOnePos{} on 98) from zero. The path $A \to B \to C$ splits it into
an augmentation leg, $B - A$, of \augprecAug{} \augprecAugCI, and a precision leg, $C - B$, of
\augprecPrec{} \augprecPrecCI; neither leg's interval excludes zero. The legs sum to the total
only because the path is a chain, and without cell~D the interaction is not identified. The
augmentation leg has the larger point estimate, \augprecAugLenses{} positives against
\augprecPrecLenses{}, but with both intervals containing zero we do not name augmentation as
the cause. With so few trainable parameters, augmentation mostly injects noise into a fixed
feature space; with the encoder unfrozen it would normally help, so none of these values says
that augmentation is harmful.

% GENERATED FILE -- do not edit by hand.
% Rebuild:  python make_all.py (generator: make_appendix_relabel_v2.py)
% Relabelled copy of tables/A4_augprec.tex; tabular number tokens asserted identical, caption/notes token changes declared.
\begin{table}[htbp]
\centering
\caption{\textbf{Augmentation and precision in the frozen comparison, three cells of a $2\times2$.} Frozen encoder (Baseline A), development fold 0, 1921 trainable parameters, 98 test positives. Cell D was never run, so the interaction between the two factors is not identified. Differences are paired class-conditional bootstrap estimates (2000 replicates, 95\% percentile interval); the one-positive resolution is $1/98$. The two legs sum to the total exactly because $A\to B\to C$ is a chain, which is arithmetic and not evidence that the effects add.}
\label{tab:augprec2}
\footnotesize
\setlength{\tabcolsep}{4pt}
\begin{tabular}{lllrr}
\toprule
Cell & Augmentation & Precision & TPR@$10^{-3}$ & AUROC \\
\midrule
A & none & fp32 & $0.2347$ & $0.9440$ \\
B & weak & fp32 & $0.1735$ & $0.9226$ \\
C & weak & fp16 & $0.1531$ & $0.9256$ \\
D & none & fp16 & \emph{not run} & \emph{not run} \\
\bottomrule
\end{tabular}

\vspace{6pt}
\footnotesize
\setlength{\tabcolsep}{4pt}
\begin{tabular}{llrlr}
\toprule
Leg & Path & $\Delta$TPR@$10^{-3}$ & 95\% interval & Excludes 0 \\
\midrule
Augmentation (held at fp32) & $B-A$ & $-0.0612$ & $[-0.1327, +0.0000]$ & no \\
Precision (held at weak augmentation) & $C-B$ & $-0.0204$ & $[-0.0510, +0.0204]$ & no \\
Total & $C-A$ & $-0.0816$ & $[-0.1429, -0.0102]$ & yes \\
\bottomrule
\end{tabular}
\end{table}

